\chapter{Sets}
\chaptercontents{A & Sets (\S2.1) \\ B & Set operations (\S2.2) \\ C & Tuples and cartesian products (\S2.2) \\ D & Chapter 2 exercises (2.E)}%
{\fE\ 2.1: 10, 27, 31, 35\\ \fE\ 2.2: 5, 24, 26, 28, 35, 36, 37, 40, 51\\ \fE\ 2.E: 6, 33, 38\\ \fR\ 2.1: 17, 18, 23, 24, 32, 36\\ \fR\ 2.2: 4, 6, 9, 25, 41, 45, 53\\ \fR\ 2.E: 4, 5, 16, 28, 32, 41; bonus 18\\ Tested in \textbf{Quiz 2} (7 Oct).}

\hsection{Sets}

\begin{key}[{Definitions 2.1.2 and 2.1.4 \textperiodcentered\ sets, the universe, set-builder notation}]
A \term{set} is a collection of elements from a specified \term{universe of discourse} $\mathcal U$; $\mathcal U$ contains every mathematical object of interest but $\mathcal U\notin\mathcal U$. Unrestricted quantifiers range over $\mathcal U$: $\forall x,\ p(x)$ means $\forall x\in\mathcal U,\ p(x)$, and
\[
\forall x\in X,\ p(x)\equiv\forall x,\ (x\in X\Rightarrow p(x)),\qquad \exists x\in X,\ p(x)\equiv\exists x,\ (x\in X\wedge p(x)).
\]
\textbf{Set-builder notation:} $a\in\{x\mid p(x)\}\Leftrightarrow a\in\mathcal U\wedge p(a)$; $\{x\in X\mid p(x)\}=\{x\mid x\in X\wedge p(x)\}$; $\{\mathrm{expr}(x)\mid p(x)\}=\{y\mid\exists x,\ y=\mathrm{expr}(x)\wedge p(x)\}$. A set written in set-builder notation need not be an element of $\mathcal U$ (e.g.\ $\{x\mid x=x\}=\mathcal U$); this is what dissolves Russell's paradox (Exercises 2.1.1, 2.1.5).
\end{key}

\begin{trybox}[{2.1.1, 2.1.3, 2.1.5 \fO}]
\textbf{2.1.1} With $S=\{X\mid X\text{ is a set}\}$ and $R=\{X\in S\mid X\notin X\}$, determine the truth value of $R\in R$.\quad \textbf{2.1.3} Why does ``$\forall x\in X$'' expand with $\Rightarrow$ but ``$\exists x\in X$'' with $\wedge$? Verify de Morgan's laws for the expanded forms.\quad \textbf{2.1.5} Repeat 2.1.1 with the updated definitions.
\end{trybox}

\hsub{Subsets}
\begin{key}[{Definition 2.1.6}]
$U$ is a \term{subset} of $X$, written $U\subseteq X$, if $\forall a\in U,\ a\in X$ (equivalently $\forall a,\ (a\in U\Rightarrow a\in X)$). $U\nsubseteq X$: not a subset. $U\subsetneq X$ (\term{proper subset}): $U\subseteq X$ and $U\ne X$. Example 2.1.11: $\N\subseteq\Z\subseteq\Q\subseteq\R\subseteq\C$.
\end{key}
\begin{strategy}[{Strategy 2.1.7 \textperiodcentered\ proving $U\subseteq X$}]
``Let $a\in U$.'' Unfold what that means. \dots\ ``So $a\in X$.'' (It is a $\forall$-proof, Strategy 1.2.10.) Example 2.1.9: $[c,d]\subseteq(a,b)$ when $a<c<d<b$, because $c\le x\le d$ gives $a<x<b$.
\end{strategy}
\begin{result}[{Proposition 2.1.12 \textperiodcentered\ transitivity}]
If $X\subseteq Y$ and $Y\subseteq Z$ then $X\subseteq Z$. \emph{Proof:} let $a\in X$; then $a\in Y$ (as $X\subseteq Y$); then $a\in Z$ (as $Y\subseteq Z$).
\end{result}

\begin{bookex}[essential]{2.1.10}
Let $a,b,c,d\in\R$ with $a<b$ and $c<d$. Prove that $[a,b)\subseteq(c,d]$ if and only if $a>c$ and $b\le d$.
\solution
($\Leftarrow$) Suppose $a>c$ and $b\le d$. Let $x\in[a,b)$, so $a\le x<b$. Then $c<a\le x$ gives $c<x$, and $x<b\le d$ gives $x\le d$. So $x\in(c,d]$. Hence $[a,b)\subseteq(c,d]$.\\
($\Rightarrow$) Suppose $[a,b)\subseteq(c,d]$. Since $a<b$, we have $a\in[a,b)$, so $a\in(c,d]$, which gives $c<a$ (and $a\le d$). For $b\le d$, suppose towards a contradiction that $b>d$. Define $x=\frac{d+b}2$. Then $d<x<b$, and $x>d\ge a$, so $x\in[a,b)$; but $x>d$ means $x\notin(c,d]$, contradicting $[a,b)\subseteq(c,d]$. Hence $b\le d$.\qed
\end{bookex}

\hsub{Set equality}
\begin{result}[{Axiom 2.1.14 \textperiodcentered\ set extensionality}]
$X=Y$ if and only if $\forall a,\ (a\in X\Leftrightarrow a\in Y)$, equivalently $X\subseteq Y$ and $Y\subseteq X$. Consequence: order and repetition in list notation do not matter.
\end{result}
\begin{strategy}[{Strategy 2.1.15 \textperiodcentered\ proof by double containment}]
($\subseteq$) Let $a\in X$ \dots\ so $a\in Y$. \quad ($\supseteq$) Let $a\in Y$ \dots\ so $a\in X$. \quad Example 2.1.16: $\{x\in\R\mid x^2\le1\}=[-1,1]$, using $(1-a)(1+a)\ge0$ for $\subseteq$ and $|a|\le1\Rightarrow a^2\le1$ for $\supseteq$.
\end{strategy}

\begin{bookex}[recommended]{2.1.17}
Prove that $\{x\in\R\mid x^2<x\}=(0,1)$.
\solution
($\subseteq$) Let $a\in\{x\in\R\mid x^2<x\}$, so $a\in\R$ and $a^2<a$, i.e.\ $a(a-1)=a^2-a<0$. A product of two real numbers is negative only if one factor is positive and the other negative. If $a>0$ and $a-1<0$ then $0<a<1$. If $a<0$ and $a-1>0$ then $a<0$ and $a>1$, impossible. So $0<a<1$, i.e.\ $a\in(0,1)$.\\
($\supseteq$) Let $a\in(0,1)$, so $0<a<1$. Multiplying $a<1$ by $a>0$ gives $a^2<a$; and $a\in\R$. So $a\in\{x\in\R\mid x^2<x\}$.\qed
\end{bookex}

\begin{bookex}[recommended]{2.1.18}
Prove by double containment that $\{0,1\}=\{1,0\}$ and $\{0,0\}=\{0\}$.
\solution
($\subseteq$) Let $a\in\{0,1\}$. Then $a=0$ or $a=1$; in either case $a$ is one of the listed elements $1,0$, so $a\in\{1,0\}$. ($\supseteq$) Let $a\in\{1,0\}$; then $a=1$ or $a=0$, so $a\in\{0,1\}$. Hence $\{0,1\}=\{1,0\}$.\\
($\subseteq$) Let $a\in\{0,0\}$. Then $a=0$ or $a=0$, so $a=0$, so $a\in\{0\}$. ($\supseteq$) Let $a\in\{0\}$; then $a=0$, so $a\in\{0,0\}$. Hence $\{0,0\}=\{0\}$.\qed
\end{bookex}

\hsub{Empty and non-empty sets}
\begin{key}[{Definition 2.1.19, Strategy 2.1.20, Theorem 2.1.25, Notation 2.1.26}]
$X$ is \term{non-empty} if it has at least one element ($\exists a,\ a\in X$); otherwise \term{empty} ($\neg\exists a,\ a\in X$). \textbf{To prove non-empty:} exhibit an element. \textbf{To prove empty:} assume there is $a\in X$ and derive a contradiction. Any two empty sets are equal (2.1.25), so there is one \term{empty set}, $\varnothing$. Examples 2.1.21--22: $\{x\in\Q\mid x^2=2\}$ and $[0]$ are empty.
\end{key}

\begin{bookex}[recommended]{2.1.23}
Let $a,b\in\R$. Prove that $[a,b]$ is empty iff $a>b$, and that $(a,b)$ is empty iff $a\ge b$.
\solution
\emph{$[a,b]$.} ($\Leftarrow$) Suppose $a>b$. If $x\in[a,b]$ then $a\le x\le b$, so $a\le b$, contradiction. So $[a,b]$ is empty. ($\Rightarrow$) By contraposition: suppose $a\le b$. Then $a\le a\le b$, so $a\in[a,b]$ and $[a,b]$ is non-empty.\\
\emph{$(a,b)$.} ($\Leftarrow$) Suppose $a\ge b$. If $x\in(a,b)$ then $a<x<b$, so $a<b$, contradiction. ($\Rightarrow$) By contraposition: suppose $a<b$. Then $\frac{a+b}2\in(a,b)$ (as in Exercise 1.2.29), so $(a,b)$ is non-empty.\qed
\end{bookex}

\begin{bookex}[recommended]{2.1.24}
Let $E$ be empty and $p(x)$ a predicate on $E$. Show that $\forall x\in E,\ p(x)$ is true and $\exists x\in E,\ p(x)$ is false. What does $\forall x\in E,\ x\ne x$ mean, and is it true?
\solution
$\forall x\in E,\ p(x)$ means $\forall x,\ (x\in E\Rightarrow p(x))$ (Definition 2.1.2). Let $x$ be arbitrary and suppose $x\in E$. Since $E$ is empty this is a contradiction, so $p(x)$ follows by the principle of explosion (Axiom 1.1.67). Hence the implication holds for every $x$, and the statement is true (``vacuously'').\\
$\exists x\in E,\ p(x)$ means $\exists x,\ (x\in E\wedge p(x))$; if it were true there would be some $x\in E$, contradicting emptiness. So it is false.\\
$\forall x\in E,\ x\ne x$ says ``every element of the empty set is different from itself''. It is true, by the first part with $p(x)$ being $x\ne x$.\qed
\end{bookex}

\begin{bookex}[essential]{2.1.27}
Is it true that $\varnothing\in X$ for all sets $X$? Is it true that $\varnothing\subseteq X$ for all sets $X$?
\solution
$\varnothing\in X$ for all $X$ is \textbf{false}: take $X=\varnothing$; it has no elements, so $\varnothing\notin\varnothing$.\\
$\varnothing\subseteq X$ for all $X$ is \textbf{true}: by Definition 2.1.6 we need $\forall a\in\varnothing,\ a\in X$, which holds by Exercise 2.1.24 (vacuously).\qed
\end{bookex}

\begin{trybox}[{2.1.28 \fO}]
\textbf{2.1.28} Why are $\varnothing$, $\{\varnothing\}$, $\{\{\varnothing\}\}$ different? Truth values of $\varnothing\in\varnothing$, $\varnothing\subseteq\varnothing$, $\varnothing\in\{\varnothing\}$, $\varnothing\in\{\{\varnothing\}\}$, $\{\varnothing\}\subseteq\{\varnothing\}$, $\{\varnothing\}\subseteq\{\{\varnothing\}\}$; find more.
\end{trybox}

\hsub{Power sets}
\begin{key}[{Definition 2.1.29, Example 2.1.34}]
$\PS(X)$ is the set of all subsets of $X$: $\quad U\in\PS(X)\ \Leftrightarrow\ U\subseteq X.$\quad $\PS(\{1,2\})=\{\varnothing,\{1\},\{2\},\{1,2\}\}$. Confusions come from mixing $\in$ (element) with $\subseteq$ (subset): $U\subseteq X$ is the same as $U\in\PS(X)$.
\end{key}

\begin{bookex}[essential]{2.1.31}
Write out the elements of $\PS(\{1,2,3\})$.
\solution
$\PS(\{1,2,3\})=\big\{\varnothing,\ \{1\},\ \{2\},\ \{3\},\ \{1,2\},\ \{1,3\},\ \{2,3\},\ \{1,2,3\}\big\}$: eight subsets, one for each choice of which of $1,2,3$ to include.\qed
\end{bookex}

\begin{bookex}[recommended]{2.1.32}
Let $X$ be a set. Show that $\varnothing\in\PS(X)$ and $X\in\PS(X)$.
\solution
By Exercise 2.1.27, $\varnothing\subseteq X$, so $\varnothing\in\PS(X)$ by Definition 2.1.29. By Example 2.1.8, $X\subseteq X$, so $X\in\PS(X)$.\qed
\end{bookex}

\begin{bookex}[essential]{2.1.35}
Let $X$ and $Y$ be sets. Prove that if $X\subseteq Y$, then $\PS(X)\subseteq\PS(Y)$.
\solution
Suppose $X\subseteq Y$. Let $U\in\PS(X)$. Then $U\subseteq X$ (Definition 2.1.29). Since $U\subseteq X$ and $X\subseteq Y$, transitivity (Proposition 2.1.12) gives $U\subseteq Y$, so $U\in\PS(Y)$. Hence $\PS(X)\subseteq\PS(Y)$ by Strategy 2.1.7.\qed
\end{bookex}

\begin{bookex}[recommended]{2.1.36}
True or false? (a) $\PS(\varnothing)\in\{\varnothing\}$; (b) $\PS(\varnothing)\in\PS(\PS(\varnothing))$; (c) $\PS(\PS(\varnothing))\in\{\varnothing,\{\varnothing,\{\varnothing\}\}\}$. Repeat with $\in$ replaced by $\subseteq$.
\solution
First compute: $\PS(\varnothing)=\{\varnothing\}$ and $\PS(\PS(\varnothing))=\PS(\{\varnothing\})=\{\varnothing,\{\varnothing\}\}$.\\
(a) False: the only element of $\{\varnothing\}$ is $\varnothing$, and $\{\varnothing\}\ne\varnothing$ (it has an element).\\
(b) True: $\PS(\varnothing)\subseteq\PS(\varnothing)$ (Example 2.1.8), so $\PS(\varnothing)\in\PS(\PS(\varnothing))$ by Definition 2.1.29.\\
(c) True: $\PS(\PS(\varnothing))=\{\varnothing,\{\varnothing\}\}$ is the second listed element.\\
With $\subseteq$: (a$'$) $\{\varnothing\}\subseteq\{\varnothing\}$: true (Example 2.1.8). (b$'$) $\{\varnothing\}\subseteq\{\varnothing,\{\varnothing\}\}$: true, its only element $\varnothing$ is listed. (c$'$) $\{\varnothing,\{\varnothing\}\}\subseteq\{\varnothing,\{\varnothing,\{\varnothing\}\}\}$: false, because $\{\varnothing\}$ is an element of the left set but not of the right one (the right set's elements are $\varnothing$ and $\{\varnothing,\{\varnothing\}\}$, and $\{\varnothing\}$ differs from both).\qed
\end{bookex}

\begin{trybox}[{2.1.33 \fO}]
\textbf{2.1.33} Write out $\PS(\varnothing)$, $\PS(\PS(\varnothing))$ and $\PS(\PS(\PS(\varnothing)))$.
\end{trybox}

\hsection{Set operations}

\hsub{Intersection}
\begin{key}[{Definitions 2.2.1, 2.2.2, 2.2.10, Notation 2.2.13}]
The \term{intersection} of a collection of sets is the set of objects that are elements of all of them. $A\cap B=\{x\mid x\in A\text{ and }x\in B\}$, so $x\in A\cap B\Leftrightarrow x\in A\wedge x\in B$. $A$ and $B$ are \term{disjoint} if $A\cap B$ is empty. $A_0\cap\dots\cap A_{n-1}=\{x\mid x\in A_i\text{ for all }i\in[n]\}$.
\end{key}

\begin{bookex}[recommended]{2.2.4}
Prove that $[0,\infty)\cap\Z=\N$.
\solution
($\subseteq$) Let $x\in[0,\infty)\cap\Z$. Then $x\in[0,\infty)$ and $x\in\Z$, so $x$ is an integer with $x\ge0$, i.e.\ a non-negative whole number, so $x\in\N$ (Definition 0.6). ($\supseteq$) Let $n\in\N$. Then $n$ is a whole number, so $n\in\Z$, and $n\ge0$, so $n\in[0,\infty)$. Hence $n\in[0,\infty)\cap\Z$.\qed
\end{bookex}

\begin{bookex}[essential]{2.2.5}
Write down the elements of $\{0,1,4,7\}\cap\{1,2,3,4,5\}$.
\solution
The elements common to both lists are $1$ and $4$: $\{0,1,4,7\}\cap\{1,2,3,4,5\}=\{1,4\}$.\qed
\end{bookex}

\begin{bookex}[recommended]{2.2.6}
Express $[-2,5)\cap[4,7)$ as a single interval.
\solution
$x$ lies in both iff $-2\le x<5$ and $4\le x<7$, iff $4\le x<5$. So $[-2,5)\cap[4,7)=[4,5)$.\qed
\end{bookex}

\begin{result}[{Proposition 2.2.8}]
$X\subseteq Y$ if and only if $X\cap Y=X$. \emph{Proof:} ($\Rightarrow$) double containment, using $X\subseteq Y$ for $\supseteq$; ($\Leftarrow$) if $a\in X=X\cap Y$ then $a\in Y$.
\end{result}

\begin{bookex}[recommended]{2.2.9}
Let $X$ be a set. Prove that $X\cap\varnothing=\varnothing$ and $X\cap X=X$.
\solution
\emph{$X\cap\varnothing=\varnothing$.} Suppose $X\cap\varnothing$ is non-empty, say $a\in X\cap\varnothing$. Then $a\in\varnothing$, contradicting that $\varnothing$ has no elements. So $X\cap\varnothing$ is empty, hence equal to $\varnothing$ (Theorem 2.1.25).\\
\emph{$X\cap X=X$.} For any $a$: $a\in X\cap X\Leftrightarrow a\in X\wedge a\in X\Leftrightarrow a\in X$. So the sets are equal by extensionality (Axiom 2.1.14).\qed
\end{bookex}

\begin{trybox}[{2.2.7, 2.2.12, 2.2.15 \fO}]
\textbf{2.2.7} For $A\subseteq\R$, express with $\cap$: (a) the positive elements of $A$; (b) the integer elements of $A$.\quad \textbf{2.2.12} Prove that $(a,b)$ and $(c,d)$ are disjoint iff $b\le c$ or $d\le a$.\quad \textbf{2.2.15} Express the set of non-negative rational elements of $A$ using $\cap$.
\end{trybox}

\hsub{Indexed families}
\begin{key}[{Definition 2.2.16, Notation 2.2.18}]
An \term{indexed family} over a set $I$ assigns a set $A_i$ to each $i\in I$: $\{A_i\mid i\in I\}$. $\displaystyle\bigcap_{i\in I}A_i=\{x\mid x\in A_i\text{ for all }i\in I\}$, so $x\in\bigcap_{i\in I}A_i\Leftrightarrow\forall i\in I,\ x\in A_i$. For $I=\{n\in\N\mid n\ge1\}$ one writes $\bigcap_{n\ge1}A_n$ or $\bigcap_{n=1}^\infty A_n$.
\end{key}
\begin{strategy}[{Example 2.2.19 \textperiodcentered\ the technique for intersections over $n\ge1$}]
To prove $\bigcap_{n\ge1}[0,1+\frac1n)=[0,1]$: ($\supseteq$) let $x\in[0,1]$ and let $n\ge1$ be arbitrary; check $x<1+\frac1n$. ($\subseteq$) let $x$ be in every $[0,1+\frac1n)$ and suppose $x>1$; choose $N=\lfloor\frac1{x-1}\rfloor+1$, which gives $N\ge1$ and $N\ge\frac1{x-1}$, hence $x\ge1+\frac1N$, contradicting $x\in[0,1+\frac1N)$. \textbf{The invention is always the same: build an index $N$ from $x$ that breaks the inequality.}
\end{strategy}

\begin{trybox}[{2.2.20 \fO}]
\textbf{2.2.20} For $n\in\Z$ let $A_n=\{x\in\R\mid x\ne n\}$. Describe $\bigcap_{n\in\Z}A_n$ in plain terms.
\end{trybox}

\hsub{Union}
\begin{key}[{Definitions 2.2.21, 2.2.22, Notations 2.2.29, 2.2.33}]
The \term{union} of a collection of sets is the set of objects that are elements of at least one of them. $A\cup B=\{x\mid x\in A\text{ or }x\in B\}$, so $x\in A\cup B\Leftrightarrow x\in A\vee x\in B$. $A_0\cup\dots\cup A_{n-1}=\{x\mid x\in A_i\text{ for some }i\in[n]\}$. $\displaystyle\bigcup_{i\in I}A_i=\{x\mid x\in A_i\text{ for some }i\in I\}$, so $x\in\bigcup_{i\in I}A_i\Leftrightarrow\exists i\in I,\ x\in A_i$. Example 2.2.23: $E\cup O=\Z$ and $E\cap O=\varnothing$ for the even and odd integers.
\end{key}

\begin{bookex}[essential]{2.2.24}
Write down the elements of $\{0,1,4,7\}\cup\{1,2,3,4,5\}$.
\solution
Everything in either list, each element once: $\{0,1,2,3,4,5,7\}$.\qed
\end{bookex}

\begin{bookex}[recommended]{2.2.25}
Express $[-2,5)\cup[4,7)$ as a single interval.
\solution
$x$ lies in the union iff $-2\le x<5$ or $4\le x<7$; since the two ranges overlap on $[4,5)$, this is iff $-2\le x<7$. So $[-2,5)\cup[4,7)=[-2,7)$.\qed
\end{bookex}

\begin{bookex}[essential]{2.2.26}
Let $X$ and $Y$ be sets. Prove that $X\subseteq Y$ if and only if $X\cup Y=Y$.
\solution
($\Rightarrow$) Suppose $X\subseteq Y$; we prove $X\cup Y=Y$ by double containment. ($\subseteq$) Let $a\in X\cup Y$. Then $a\in X$ or $a\in Y$. If $a\in X$ then $a\in Y$ since $X\subseteq Y$; if $a\in Y$ then $a\in Y$. In both cases $a\in Y$. ($\supseteq$) Let $a\in Y$. Then $a\in X\cup Y$ by definition of union.\\
($\Leftarrow$) Suppose $X\cup Y=Y$. Let $a\in X$. Then $a\in X\cup Y$ by definition of union, so $a\in Y$ since $X\cup Y=Y$. Hence $X\subseteq Y$.\qed
\end{bookex}

\begin{result}[{Example 2.2.27 \textperiodcentered\ distributivity of $\cap$ over $\cup$}]
$X\cap(Y\cup Z)=(X\cap Y)\cup(X\cap Z)$, by double containment; the $\subseteq$ direction splits into cases on $x\in Y\cup Z$.
\end{result}

\begin{bookex}[essential]{2.2.28}
Let $X,Y,Z$ be sets. Prove that $X\cup(Y\cap Z)=(X\cup Y)\cap(X\cup Z)$.
\solution
($\subseteq$) Let $x\in X\cup(Y\cap Z)$. Then $x\in X$ or $x\in Y\cap Z$. If $x\in X$, then $x\in X\cup Y$ and $x\in X\cup Z$, so $x\in(X\cup Y)\cap(X\cup Z)$. If $x\in Y\cap Z$, then $x\in Y$ and $x\in Z$, so $x\in X\cup Y$ and $x\in X\cup Z$, and again $x\in(X\cup Y)\cap(X\cup Z)$.\\
($\supseteq$) Let $x\in(X\cup Y)\cap(X\cup Z)$, so $x\in X\cup Y$ and $x\in X\cup Z$. By the law of excluded middle, $x\in X$ or $x\notin X$. If $x\in X$, then $x\in X\cup(Y\cap Z)$. If $x\notin X$, then from $x\in X\cup Y$ we get $x\in Y$, and from $x\in X\cup Z$ we get $x\in Z$; so $x\in Y\cap Z$, and $x\in X\cup(Y\cap Z)$.\qed
\end{bookex}

\begin{trybox}[{2.2.31, 2.2.32 \fO}]
\textbf{2.2.31} Write $X=\{x\in\R\mid x^4-5x^2+4>0\}$ as a union of intervals.\quad \textbf{2.2.32} State and prove the distributivity laws for finitary unions and intersections.
\end{trybox}

\begin{bookex}[essential]{2.2.35}
Express $\bigcup_{n\ge1}[0,1-\frac1n]$ as an interval.
\solution
We claim $\bigcup_{n\ge1}[0,1-\frac1n]=[0,1)$. ($\subseteq$) Let $x$ be in the union; take $n\ge1$ with $x\in[0,1-\frac1n]$. Then $0\le x\le1-\frac1n<1$, so $x\in[0,1)$. ($\supseteq$) Let $x\in[0,1)$, so $0\le x<1$ and $1-x>0$. Define $N=\lceil\frac1{1-x}\rceil$; since $\frac1{1-x}\ge1$ we have $N\ge1$, and $N\ge\frac1{1-x}$ gives $\frac1N\le1-x$, i.e.\ $x\le1-\frac1N$. So $x\in[0,1-\frac1N]$ and hence $x$ is in the union.\qed
\end{bookex}

\begin{bookex}[essential]{2.2.36}
Prove that $\bigcap_{n\in\N}[n]=\varnothing$ and $\bigcup_{n\in\N}[n]=\N$.
\solution
\emph{Intersection.} Suppose $x\in\bigcap_{n\in\N}[n]$. Then $x\in[n]$ for all $n\in\N$, in particular $x\in[0]$. But $[0]=\{k\in\N\mid k<0\}$ is empty (Example 2.1.22): contradiction. So the intersection is empty.\\
\emph{Union.} ($\subseteq$) If $x\in[n]$ for some $n$, then $x\in\N$ by definition of $[n]$. ($\supseteq$) Let $k\in\N$. Then $k<k+1$ and $k+1\in\N$, so $k\in[k+1]$, hence $k\in\bigcup_{n\in\N}[n]$.\qed
\end{bookex}

\begin{bookex}[essential]{2.2.37}
Find a family $\{X_n\mid n\in\N\}$ with (i) $\bigcup_{n\in\N}X_n=\N$; (ii) $\bigcap_{n\in\N}X_n=\varnothing$; (iii) $X_i\cap X_j\ne\varnothing$ for all $i,j\in\N$.
\solution
Take $X_n=\{k\in\N\mid k\ne n\}$ (``$\N$ with $n$ removed''). (i) Each $X_n\subseteq\N$; and any $k\in\N$ lies in $X_{k+1}$ since $k\ne k+1$. (ii) If $k$ were in every $X_n$ then in particular $k\in X_k$, i.e.\ $k\ne k$: contradiction. (iii) Given $i,j$, the number $i+j+1$ differs from both $i$ and $j$, so $i+j+1\in X_i\cap X_j$.\qed
\end{bookex}

\hsub{Relative complement}
\begin{key}[{Definition 2.2.38}]
$X\setminus A=\{x\in X\mid x\notin A\}$, the \term{relative complement} of $A$ in $X$ (``$X$ minus $A$''); $x\in X\setminus A\Leftrightarrow x\in X\wedge x\notin A$. The book uses no absolute complement $A^c$. Example 2.2.39: $\Z\setminus E$ is the set of odd integers.
\end{key}

\begin{bookex}[essential]{2.2.40}
Write down the elements of $\{0,1,4,7\}\setminus\{1,2,3,4,5\}$.
\solution
Keep the elements of the first set that are not in the second: $\{0,7\}$.\qed
\end{bookex}

\begin{bookex}[recommended]{2.2.41}
Express $[-2,5)\setminus[4,7)$ and $[4,7)\setminus[-2,5)$ as intervals.
\solution
$x\in[-2,5)\setminus[4,7)$ iff $-2\le x<5$ and not $4\le x<7$, iff $-2\le x<4$: the set is $[-2,4)$.\\
$x\in[4,7)\setminus[-2,5)$ iff $4\le x<7$ and not $-2\le x<5$, iff $5\le x<7$: the set is $[5,7)$.\qed
\end{bookex}

\begin{trybox}[{2.2.42, 2.2.43 \fO}]
\textbf{2.2.42} $A$ = the reals that are not integers: (a) write $A$ as a relative complement; (b) as a union of intervals.\quad \textbf{2.2.43} Prove $X\setminus(X\setminus A)=X\cap A$, and deduce that $A\subseteq X$ iff $X\setminus(X\setminus A)=A$.
\end{trybox}

\hsub{Sets versus logic}
\begin{key}[{The dictionary (page 92)}]
With $p$: $a\in X$, $q$: $a\in Y$, $r(i)$: $a\in X_i$:
\[
\begin{array}{ll@{\qquad}ll}
a\notin X & \neg p & a\in\bigcap_{i\in I}X_i & \forall i\in I,\ r(i)\\
a\in X\cap Y & p\wedge q & a\in\bigcup_{i\in I}X_i & \exists i\in I,\ r(i)\\
a\in X\cup Y & p\vee q & a\in X\setminus Y & p\wedge(\neg q)
\end{array}
\]
Every set identity in this section is a logical law in disguise; the proof unfolds membership, applies the law, folds back.
\end{key}
\begin{result}[{Theorem 2.2.44 \textperiodcentered\ de Morgan's laws for sets}]
(a) $X\setminus(A\cup B)=(X\setminus A)\cap(X\setminus B)$;\ (b) $X\setminus(A\cap B)=(X\setminus A)\cup(X\setminus B)$;\ (c) $X\setminus\bigcup_{i\in I}A_i=\bigcap_{i\in I}(X\setminus A_i)$;\ (d) $X\setminus\bigcap_{i\in I}A_i=\bigcup_{i\in I}(X\setminus A_i)$. The book proves (a) by double containment, using Theorem 1.3.24 to turn $x\notin A\cup B$ into $x\notin A$ and $x\notin B$.
\end{result}

\begin{bookex}[recommended]{2.2.45}
Complete the proof of de Morgan's laws for sets: parts (b), (c), (d).
\solution
\emph{(b)} ($\subseteq$) Let $x\in X\setminus(A\cap B)$: $x\in X$ and $x\notin A\cap B$, i.e.\ $\neg(x\in A\wedge x\in B)$. By de Morgan (Theorem 1.3.24(a)), $x\notin A$ or $x\notin B$. If $x\notin A$ then $x\in X\setminus A$; if $x\notin B$ then $x\in X\setminus B$. Either way $x\in(X\setminus A)\cup(X\setminus B)$. ($\supseteq$) Let $x\in(X\setminus A)\cup(X\setminus B)$. If $x\in X\setminus A$: $x\in X$ and $x\notin A$, so $x\notin A\cap B$ (an element of $A\cap B$ would be in $A$), hence $x\in X\setminus(A\cap B)$. The case $x\in X\setminus B$ is the same.\\
\emph{(c)} For any $x$: $x\in X\setminus\bigcup_iA_i\Leftrightarrow x\in X\wedge\neg(\exists i\in I,\ x\in A_i)\Leftrightarrow x\in X\wedge\forall i\in I,\ x\notin A_i$ \reason{Theorem 1.3.29(b)} $\Leftrightarrow\forall i\in I,\ (x\in X\wedge x\notin A_i)$ \reason{$x\in X$ does not depend on $i$; see note} $\Leftrightarrow\forall i\in I,\ x\in X\setminus A_i\Leftrightarrow x\in\bigcap_i(X\setminus A_i)$. \emph{Note:} the third step needs $I$ non-empty for the direction $\Leftarrow$; for $I=\varnothing$ one reads the conventions of the book for empty intersections.\\
\emph{(d)} For any $x$: $x\in X\setminus\bigcap_iA_i\Leftrightarrow x\in X\wedge\neg(\forall i\in I,\ x\in A_i)\Leftrightarrow x\in X\wedge\exists i\in I,\ x\notin A_i$ \reason{Theorem 1.3.29(a)} $\Leftrightarrow\exists i\in I,\ (x\in X\wedge x\notin A_i)\Leftrightarrow\exists i\in I,\ x\in X\setminus A_i\Leftrightarrow x\in\bigcup_i(X\setminus A_i)$.\qed
\end{bookex}

\begin{warn}
\begin{itemize}
\item From $x\in A\cup B$ you get a \emph{disjunction}: do cases. From $x\in A\cap B$ you get two facts. From $x\notin A\cup B$ you get $x\notin A$ and $x\notin B$ (de Morgan), not just one of them.
\item $X\setminus A=X\setminus B$ does not give $A=B$ (aside on page 91): the complement operation does not cancel like subtraction.
\item An ``express as an interval'' question still deserves the one-line justification ``$x$ is in the set iff \dots''.
\end{itemize}
\end{warn}
\hsection{Tuples and cartesian products}

\begin{key}[{Definition 2.2.46 \textperiodcentered\ ordered lists (tuples)}]
A (finite) \term{ordered list} is an expression $(x_0,x_1,\dots,x_{k-1})$ with $k\in\N$; $k$ is its \term{length} and $i\in[k]$ is the \term{position} of $x_i$. Two lists are equal exactly when they have the same length and the same objects in all positions:
\[
(x_0,\dots,x_{k-1})=(y_0,\dots,y_{\ell-1})\ \Leftrightarrow\ k=\ell\ \text{and}\ x_i=y_i\ \text{for all } i\in[k].
\]
A list of length $k$ is a \term{$k$-tuple}; $()$ is the empty list; the 1-tuple $(x)$ is identified with $x$; 2-, 3-, 4-tuples are \term{ordered pairs, triples, quadruples}. Example 2.2.47: $(1,2,3)\ne(1,2,3,3,3)$ (different lengths) and $(1,2,3)\ne(3,2,1)$ (different objects in position $0$), although $\{1,2,3\}=\{3,2,1\}=\{1,2,3,3,3\}$.
\end{key}

\begin{bookex}{2.2.48}
Write down all ordered lists that have all of the elements of $\{1,2,3\}$ as terms and have no repeated terms.
\solution
Such a list has length $3$ (three distinct terms, each of $1,2,3$ appearing). Choosing the term in position $0$, then position $1$, leaves the last term forced:
\[
(1,2,3),\ (1,3,2),\ (2,1,3),\ (2,3,1),\ (3,1,2),\ (3,2,1).
\]
These six lists are pairwise different by the equality rule of Definition 2.2.46 (they differ in some position).\qed
\end{bookex}

\begin{key}[{Definition 2.2.49 and Notation 2.2.54 \textperiodcentered\ cartesian product, $A^k$}]
For $k\in\N$ and sets $A_0,\dots,A_{k-1}$,
\[
\prod_{i\in[k]}A_i=A_0\times A_1\times\dots\times A_{k-1}=\{(x_0,x_1,\dots,x_{k-1})\mid x_i\in A_i\text{ for all }i\in[k]\}.
\]
So $(x,y)\in X\times Y$ means exactly $x\in X$ and $y\in Y$. The \term{$k$-fold product} is $A^k=A\times\dots\times A$ ($k$ times), the set of all $k$-tuples of elements of $A$; $A^0=\{()\}$ and $A^1=A$. Example 2.2.50: $\R\times\R=\{(x,y)\mid x,y\in\R\}$ is the plane, $\R\times\{0\}$ is the $x$-axis, and the graph of $y=x^2$ is $\{(x,x^2)\mid x\in\R\}\subseteq\R\times\R$.
\end{key}

\begin{bookex}[essential]{2.2.51}
Write down the elements of the set $\{1,2\}\times\{3,4,5\}$.
\solution
By Definition 2.2.49 the elements are the ordered pairs $(x,y)$ with $x\in\{1,2\}$ and $y\in\{3,4,5\}$:
\[
\{1,2\}\times\{3,4,5\}=\{(1,3),(1,4),(1,5),(2,3),(2,4),(2,5)\}.
\]
(Two choices for the first term, three for the second: $2\cdot3=6$ pairs.)\qed
\end{bookex}

\begin{bookex}{2.2.52}
Let $X$ be a set. Prove that $X\times\varnothing=\varnothing$.
\solution
By Strategy 2.1.20 we prove that $X\times\varnothing$ has no elements. Suppose $t\in X\times\varnothing$. Then $t=(x,y)$ with $x\in X$ and $y\in\varnothing$ \reason{Definition 2.2.49}. But $\varnothing$ has no elements, so $y\in\varnothing$ is a contradiction. Hence $X\times\varnothing$ is empty, and by Theorem 2.1.25 it equals $\varnothing$.\qed
\end{bookex}

\begin{bookex}[recommended]{2.2.53}
Let $X$ and $Y$ be sets. Under what conditions is it true that $X\times Y=Y\times X$?
\solution
\textbf{Claim.} $X\times Y=Y\times X$ if and only if $X=Y$ or $X=\varnothing$ or $Y=\varnothing$.\\
($\Leftarrow$) If $X=Y$ the two sets are literally the same. If $X=\varnothing$ then $X\times Y=\varnothing\times Y=\varnothing$ (as in Exercise 2.2.52: a pair in it would have a first term in $\varnothing$) and $Y\times X=Y\times\varnothing=\varnothing$ by Exercise 2.2.52, so both sides equal $\varnothing$. The case $Y=\varnothing$ is the same.\\
($\Rightarrow$) Suppose $X\times Y=Y\times X$, and suppose moreover that $X\ne\varnothing$ and $Y\ne\varnothing$; we prove $X=Y$ by double containment. Let $x\in X$. Since $Y$ is non-empty, fix some $y\in Y$. Then $(x,y)\in X\times Y=Y\times X$, so $x\in Y$ \reason{first term of an element of $Y\times X$}. Hence $X\subseteq Y$. Symmetrically, let $y\in Y$ and fix $x\in X$; then $(y,x)\in Y\times X=X\times Y$, so $y\in X$. Hence $Y\subseteq X$, and $X=Y$.\qed
\end{bookex}

\begin{trybox}[{2.2.55 \fO}]
\textbf{2.2.55} Are $\R^3$, $\R^2\times\R$, $\R\times\R^2$, $\R\times\R\times\R$, $(\R\times\R)\times\R$, $\R\times(\R\times\R)$ all equal? Which pairs are equal, and in what other sense are they all equivalent? \hint{write down a typical element of each: is it a triple, or a pair one of whose terms is a pair?}
\end{trybox}

\hsection{Chapter 2 exercises}

\hsub{Sets, subsets and power sets (2.E.1--2.E.7)}
\begin{strategy}[{What the chapter exercises test}]
Everything in 2.E is done with four moves from \S2.1--\S2.2: \textbf{(1)} unfold membership (``$a\in X\setminus V$ means $a\in X$ and $a\notin V$''); \textbf{(2)} prove $\subseteq$ by ``let $a\in\dots$'' (Strategy 2.1.7) and $=$ by double containment (Strategy 2.1.15); \textbf{(3)} prove a set is empty by assuming it has an element and deriving a contradiction (Strategy 2.1.20); \textbf{(4)} to refute a general claim, give a concrete counterexample, usually built from $\varnothing$, $\{0\}$, $\{1\}$, $\{0,1\}$. Remember $\PS(X)=\{U\mid U\subseteq X\}$, so $U\in\PS(X)\Leftrightarrow U\subseteq X$ (Definition 2.1.29), and $\varnothing\subseteq X$ always (Exercise 2.1.27).
\end{strategy}

\begin{bookex}[recommended]{2.E.4}
Let $X$ be a set and let $U,V\subseteq X$. Prove that $U$ and $V$ are disjoint if and only if $U\subseteq X\setminus V$.
\solution
($\Rightarrow$) Suppose $U\cap V=\varnothing$. Let $a\in U$. Then $a\in X$ since $U\subseteq X$. If $a\in V$ we would have $a\in U\cap V=\varnothing$, a contradiction; so $a\notin V$. Hence $a\in X\setminus V$ \reason{Definition 2.2.38}, and $U\subseteq X\setminus V$.\\
($\Leftarrow$) Suppose $U\subseteq X\setminus V$. Suppose towards a contradiction that $a\in U\cap V$. Then $a\in U$, so $a\in X\setminus V$, so $a\notin V$; but also $a\in V$. Contradiction. So $U\cap V$ is empty: $U$ and $V$ are disjoint.\qed
\end{bookex}

\begin{bookex}[recommended]{2.E.5}
For each statement, determine whether it is true for all sets $A$ and $X$, and prove your claim.\\
(a) If $X\setminus A=\varnothing$ then $X=A$.\quad (b) If $X\setminus A=X$ then $A=\varnothing$.\quad (c) If $X\setminus A=A$ then $A=\varnothing$.\quad (d) $X\setminus(X\setminus A)=A$.
\solution
\emph{(a) False.} Take $X=\varnothing$ and $A=\{0\}$. Then $X\setminus A=\varnothing$ but $X\ne A$. (In general $X\setminus A=\varnothing$ says only that $X\subseteq A$.)\\
\emph{(b) False.} Take $X=\varnothing$ and $A=\{0\}$. Then $X\setminus A=\varnothing=X$ but $A\ne\varnothing$. (In general $X\setminus A=X$ says only that $X$ and $A$ are disjoint.)\\
\emph{(c) True.} Suppose $X\setminus A=A$ and suppose towards a contradiction that $a\in A$. Then $a\in X\setminus A$, so $a\notin A$. Contradiction. Hence $A$ is empty.\\
\emph{(d) False.} Take $X=\varnothing$ and $A=\{0\}$: $X\setminus(X\setminus A)=\varnothing\setminus\varnothing=\varnothing\ne A$. (By Exercise 2.2.43, $X\setminus(X\setminus A)=X\cap A$, which equals $A$ exactly when $A\subseteq X$.)\qed
\end{bookex}

\begin{bookex}[essential]{2.E.6}
For each statement, determine whether it is true for all sets $X,Y$, false for all sets $X,Y$, or true for some and false for others.\\
(a) $\PS(X\cup Y)=\PS(X)\cup\PS(Y)$\quad (b) $\PS(X\cap Y)=\PS(X)\cap\PS(Y)$\quad (c) $\PS(X\times Y)=\PS(X)\times\PS(Y)$\quad (d) $\PS(X\setminus Y)=\PS(X)\setminus\PS(Y)$
\solution
\emph{(a) Sometimes.} True for $X=Y=\varnothing$: both sides are $\PS(\varnothing)=\{\varnothing\}$. False for $X=\{0\}$, $Y=\{1\}$: $\{0,1\}\in\PS(X\cup Y)$, but $\{0,1\}\nsubseteq X$ and $\{0,1\}\nsubseteq Y$, so $\{0,1\}\notin\PS(X)\cup\PS(Y)$. (The inclusion $\supseteq$ always holds: a subset of $X$ or of $Y$ is a subset of $X\cup Y$.)\\
\emph{(b) Always.} For any $U$: $U\in\PS(X\cap Y)\Leftrightarrow U\subseteq X\cap Y\Leftrightarrow(U\subseteq X\wedge U\subseteq Y)\Leftrightarrow U\in\PS(X)\wedge U\in\PS(Y)\Leftrightarrow U\in\PS(X)\cap\PS(Y)$. The middle step: if $U\subseteq X\cap Y$ then each $a\in U$ lies in $X\cap Y$, hence in $X$ and in $Y$; conversely if $U\subseteq X$ and $U\subseteq Y$ then each $a\in U$ lies in $X$ and in $Y$, hence in $X\cap Y$.\\
\emph{(c) Never.} The elements of $\PS(X\times Y)$ are \emph{sets} (subsets of $X\times Y$), whereas the elements of $\PS(X)\times\PS(Y)$ are \emph{ordered pairs} $(U,V)$. In particular $\varnothing\in\PS(X\times Y)$ (Exercise 2.1.27), but $\varnothing$ is not an ordered pair, so $\varnothing\notin\PS(X)\times\PS(Y)$. Hence the two sets are never equal.\\
\emph{(d) Never.} $\varnothing\in\PS(X\setminus Y)$ always. But $\varnothing\in\PS(Y)$, so $\varnothing\notin\PS(X)\setminus\PS(Y)$. Hence the sets are never equal.\qed
\end{bookex}

\begin{trybox}[{2.E.1, 2.E.2, 2.E.3, 2.E.7 \fO}]
\textbf{2.E.1} Express: (a) $\{n\in\Z\mid n^2<20\}$ in list notation; (b) $\{4k+3\mid k\in\N\}$ in implied list notation; (c) the odd multiples of six in set-builder notation; (d) $\{1,2,5,10,17,\dots,n^2+1,\dots\}$ in set-builder notation.\quad
\textbf{2.E.2} Find sets $X_n$ ($n\in\N$) with $X_{n+1}\subsetneq X_n$ for all $n$. Can any $X_n$ be empty?\quad
\textbf{2.E.3} Express $\PS(\{\varnothing,\{\varnothing,\{\varnothing\}\}\})$ in list notation.\quad
\textbf{2.E.7} $F$ is a set of sets with $\forall A\in F,\ \forall x\in A,\ x\in F$. Prove $F\subseteq\PS(F)$.
\end{trybox}

\hsub{Symmetric difference (2.E.8--2.E.14)}
\begin{key}[{Definition 2.E.8 \textperiodcentered\ symmetric difference}]
$X\symdiff Y=\{a\mid a\in X\text{ or }a\in Y\text{ but not both}\}$. In the dictionary of page 92: $a\in X\symdiff Y$ is $(p\vee q)\wedge\neg(p\wedge q)$ where $p$: $a\in X$, $q$: $a\in Y$. The book's hint for 2.E.13--14: rather than long chains of equations, use double containment with cases, or ``a cunning use of truth tables'' (a truth table in $p,q,r$ decides any identity built from $\cap,\cup,\setminus,\symdiff$).
\end{key}

\begin{bookex}{2.E.9}
Prove that $X\symdiff Y=(X\setminus Y)\cup(Y\setminus X)=(X\cup Y)\setminus(X\cap Y)$ for all sets $X$ and $Y$.
\solution
Let $a$ be arbitrary and write $p$ for $a\in X$ and $q$ for $a\in Y$. Then $a\in X\symdiff Y$ is $(p\vee q)\wedge\neg(p\wedge q)$; $a\in(X\setminus Y)\cup(Y\setminus X)$ is $(p\wedge\neg q)\vee(q\wedge\neg p)$; and $a\in(X\cup Y)\setminus(X\cap Y)$ is $(p\vee q)\wedge\neg(p\wedge q)$, literally the same formula as the first. For the middle one:
\[
\begin{array}{cc|c|c}
p&q&(p\vee q)\wedge\neg(p\wedge q)&(p\wedge\neg q)\vee(q\wedge\neg p)\\\hline
\checkmark&\checkmark&\times&\times\\
\checkmark&\times&\checkmark&\checkmark\\
\times&\checkmark&\checkmark&\checkmark\\
\times&\times&\times&\times
\end{array}
\]
The columns agree, so the three membership conditions are equivalent for every $a$; by set extensionality (Axiom 2.1.14) the three sets are equal.\qed
\end{bookex}

\begin{trybox}[{2.E.10--2.E.14 \fO}]
\textbf{2.E.10} $X\symdiff X=\varnothing$ and $X\symdiff\varnothing=X$.\quad \textbf{2.E.11} $X=Y\Leftrightarrow X\symdiff Y=\varnothing$.\quad \textbf{2.E.12} $X,Y$ disjoint $\Leftrightarrow X\symdiff Y=X\cup Y$.\quad \textbf{2.E.13} $X\symdiff(Y\symdiff Z)=(X\symdiff Y)\symdiff Z$.\quad \textbf{2.E.14} $X\cap(Y\symdiff Z)=(X\cap Y)\symdiff(X\cap Z)$.
\end{trybox}

\hsub{Open subsets of $\R$ (2.E.15--2.E.20)}
\begin{key}[{Definition 2.E.15 \textperiodcentered\ open set}]
$U\subseteq\R$ is \term{open} if for all $a\in U$ there exists $\delta>0$ such that $(a-\delta,a+\delta)\subseteq U$. \textbf{To prove $U$ open:} ``Let $a\in U$. Put $\delta=\dots$ \reason{in terms of $a$}; then $\delta>0$, and if $a-\delta<x<a+\delta$ then \dots\ so $x\in U$.'' \textbf{To prove $U$ not open:} negate ($\exists a\in U,\ \forall\delta>0,\ (a-\delta,a+\delta)\nsubseteq U$): ``Take $a=\dots\in U$. Let $\delta>0$. Then $x=\dots\in(a-\delta,a+\delta)$ but $x\notin U$.''
\end{key}

\begin{bookex}[recommended]{2.E.16}
For each subset of $\R$, determine (with proof) whether it is open: (a) $\varnothing$; (b) $(0,1)$; (c) $(0,1]$; (d) $\Z$; (e) $\R\setminus\Z$; (f) $\Q$.
\solution
\emph{(a) Open.} The condition ``for all $a\in\varnothing$, \dots'' is vacuously true (Exercise 2.1.24).\\
\emph{(b) Open.} Let $a\in(0,1)$ and put $\delta=\min(a,1-a)$; then $\delta>0$ since $0<a<1$. If $x\in(a-\delta,a+\delta)$ then $x>a-\delta\ge a-a=0$ and $x<a+\delta\le a+(1-a)=1$, so $x\in(0,1)$. Hence $(a-\delta,a+\delta)\subseteq(0,1)$.\\
\emph{(c) Not open.} Take $a=1\in(0,1]$ and let $\delta>0$. Then $x=1+\frac\delta2$ satisfies $1-\delta<x<1+\delta$, but $x>1$, so $x\notin(0,1]$. Hence no $\delta$ works for $a=1$.\\
\emph{(d) Not open.} Take $a=0\in\Z$ and let $\delta>0$. Put $x=\frac12\min(\delta,1)$. Then $0<x<\delta$, so $x\in(-\delta,\delta)$, and $0<x\le\frac12$, so $x$ is not an integer.\\
\emph{(e) Open.} Let $a\in\R\setminus\Z$ and let $n=\lfloor a\rfloor$, the greatest integer $\le a$; since $a\notin\Z$ we have $n<a<n+1$. Put $\delta=\min(a-n,\ n+1-a)>0$. If $x\in(a-\delta,a+\delta)$ then $n\le a-\delta<x<a+\delta\le n+1$, so $x$ lies strictly between the consecutive integers $n$ and $n+1$ and is not an integer: $x\in\R\setminus\Z$.\\
\emph{(f) Not open.} Take $a=0\in\Q$ and let $\delta>0$. Choose $n\in\N$ with $n>\sqrt2/\delta$ (the natural numbers are unbounded in $\R$) and put $x=\sqrt2/n$. Then $0<x<\delta$, so $x\in(-\delta,\delta)$; but $x\notin\Q$, for if $\sqrt2/n=r\in\Q$ then $\sqrt2=nr\in\Q$, contradicting Assumption 0.49.\qed
\end{bookex}

\begin{bookex}[bonus]{2.E.18}
(a) Let $n\ge1$ and let $U_1,\dots,U_n$ be open subsets of $\R$. Prove that $U_1\cap\dots\cap U_n$ is open. (b) Prove that $(0,1+\frac1n)$ is open for all $n\ge1$, but that $\bigcap_{n\ge1}(0,1+\frac1n)$ is not open.
\solution
\emph{(a)} Let $a\in U_1\cap\dots\cap U_n$. For each $i$ we have $a\in U_i$, and $U_i$ is open, so there is $\delta_i>0$ with $(a-\delta_i,a+\delta_i)\subseteq U_i$. Put $\delta=\min(\delta_1,\dots,\delta_n)$; then $\delta>0$ (a minimum of finitely many positive numbers). If $x\in(a-\delta,a+\delta)$ then, for each $i$, $a-\delta_i\le a-\delta<x<a+\delta\le a+\delta_i$, so $x\in U_i$. Hence $x\in U_1\cap\dots\cap U_n$, and $(a-\delta,a+\delta)\subseteq U_1\cap\dots\cap U_n$.\\
\emph{(b)} Fix $n\ge1$. Let $a\in(0,1+\frac1n)$ and put $\delta=\min(a,\ 1+\frac1n-a)>0$. If $a-\delta<x<a+\delta$ then $x>0$ and $x<1+\frac1n$, so $(a-\delta,a+\delta)\subseteq(0,1+\frac1n)$: the interval is open.\\
Now $\bigcap_{n\ge1}(0,1+\frac1n)=(0,1]$. ($\supseteq$) If $0<x\le1$ then $0<x<1+\frac1n$ for every $n\ge1$. ($\subseteq$) Let $x$ be in every $(0,1+\frac1n)$; then $x>0$. If $x>1$, choose $n\ge1$ with $n>\frac1{x-1}$; then $\frac1n<x-1$, i.e.\ $1+\frac1n<x$, so $x\notin(0,1+\frac1n)$, a contradiction. Hence $x\le1$. Finally $(0,1]$ is not open by Exercise 2.E.16(c). So an intersection of infinitely many open sets need not be open: the proof of (a) breaks down because infinitely many $\delta_i$ need not have a positive minimum.\qed
\end{bookex}

\begin{trybox}[{2.E.17, 2.E.19, 2.E.20 \fO}]
\textbf{2.E.17} $U$ is open iff for all $a\in U$ there are $u<a<v$ with $(u,v)\subseteq U$.\quad \textbf{2.E.19} $U$ is open iff $U=\bigcup_{i\in I}(a_i,b_i)$ for some family of open intervals.\quad \textbf{2.E.20} Families $\{A_i\}$, $\{B_j\}$ with (i) $\forall i\,\exists j\ge i,\ B_j\subseteq A_i$ and (ii) $\forall j\,\exists i\ge j,\ A_i\subseteq B_j$: prove $\bigcap_iA_i=\bigcap_jB_j$.
\end{trybox}

\hsub{True--false and always--sometimes--never (2.E.21--2.E.42)}
\begin{strategy}[{How to answer these}]
\textbf{True/false:} decide, then prove the statement or exhibit a counterexample. \textbf{Always:} prove the conclusion from the hypotheses for arbitrary $X,Y,\dots$. \textbf{Never:} prove the \emph{negation} of the conclusion from the hypotheses. \textbf{Sometimes:} give one concrete example satisfying the hypotheses where the conclusion holds and one where it fails. Tiny sets ($\varnothing$, $\{0\}$, $\{1\}$, $\{0,1\}$, $\{\varnothing\}$) settle almost every ``sometimes''.
\end{strategy}

\begin{bookex}[recommended]{2.E.28}
Let $a,b,c$ be arbitrary objects. Then $\{a,b\}=\{a,c\}$.
\solution
\emph{Sometimes.} If $b=c$ (for instance $a=b=c=0$) the two sets are the same set, so the conclusion holds. If $a=0$, $b=1$, $c=2$ then $1\in\{a,b\}$ but $1\notin\{a,c\}=\{0,2\}$, so $\{a,b\}\ne\{a,c\}$ by Axiom 2.1.14.\qed
\end{bookex}

\begin{bookex}[recommended]{2.E.32}
Let $X$ and $Y$ be sets with $X\cap Y=\varnothing$. Then $\forall a,\ a\in X\Rightarrow a\notin Y$.
\solution
\emph{Always.} Let $a$ be arbitrary and suppose $a\in X$. If $a\in Y$, then $a\in X\cap Y=\varnothing$, which is impossible. Hence $a\notin Y$. (This is just the unfolding of ``$X\cap Y$ is empty''.)\qed
\end{bookex}

\begin{bookex}[essential]{2.E.33}
Let $X$ and $Y$ be sets with $X\ne Y$. Then $\exists a,\ a\in X\wedge a\notin Y$.
\solution
\emph{Sometimes.} Take $X=\{0\}$ and $Y=\varnothing$: $X\ne Y$, and $a=0$ satisfies $a\in X$ and $a\notin Y$. Now take $X=\varnothing$ and $Y=\{0\}$: $X\ne Y$, but there is no $a$ with $a\in X$ at all, so the conclusion is false.\\
\emph{Why:} $X\ne Y$ means, by Axiom 2.1.14, that $\neg\forall a,\ (a\in X\Leftrightarrow a\in Y)$, i.e.\ some $a$ lies in exactly one of $X$, $Y$: either $a\in X\setminus Y$ or $a\in Y\setminus X$. The hypothesis does not tell you which, so the one-sided conclusion can fail.\qed
\end{bookex}

\begin{bookex}[essential]{2.E.38}
Let $X$ and $Y$ be sets such that $\PS(X)=\PS(Y)$. Then $X=Y$.
\solution
\emph{Always.} Since $X\subseteq X$ we have $X\in\PS(X)=\PS(Y)$, so $X\subseteq Y$. Likewise $Y\in\PS(Y)=\PS(X)$, so $Y\subseteq X$. Hence $X=Y$ by double containment.\qed
\end{bookex}

\begin{bookex}[recommended]{2.E.41}
Let $X$ and $Y$ be sets. Then $X\cup Y=X\cap Y$.
\solution
\emph{Sometimes.} For $X=Y$ it holds: $X\cup X=X=X\cap X$ (Exercise 2.2.26 with $X\subseteq X$, and Exercise 2.2.9). For $X=\{0\}$, $Y=\varnothing$ it fails: $X\cup Y=\{0\}\ne\varnothing=X\cap Y$.\\
In fact $X\cup Y=X\cap Y$ if and only if $X=Y$: if $X\cup Y=X\cap Y$ and $a\in X$ then $a\in X\cup Y=X\cap Y$, so $a\in Y$; symmetrically $Y\subseteq X$.\qed
\end{bookex}

\begin{trybox}[{2.E.21--2.E.27, 2.E.29--2.E.31, 2.E.34--2.E.37, 2.E.39, 2.E.40, 2.E.42 \fO}]
\textbf{True/false:} \textbf{2.E.21} $\neg\forall x,\ x\in E$ implies $E=\varnothing$.\ \textbf{2.E.22} $\forall x,\ \neg x\in E$ implies $E=\varnothing$.\ \textbf{2.E.23} $\{1,2,3\}=\{1,2,1,3,2,1\}$.\ \textbf{2.E.24} $\varnothing\in\varnothing$.\ \textbf{2.E.25} $\varnothing\in\{\varnothing\}$.\ \textbf{2.E.26} $\varnothing\in\{\{\varnothing\}\}$.\\
\textbf{Always/sometimes/never:} \textbf{2.E.27} $\{a,b\}=\{b,a\}$.\ \textbf{2.E.29} some $a$ with $a\in X\Leftrightarrow a\in Y$; then $X=Y$.\ \textbf{2.E.30} some $a\in X$ with $a\notin Y$; then $X=Y$.\ \textbf{2.E.31} $X\ne Y$; then $\forall a,\ a\in X\Rightarrow a\notin Y$.\ \textbf{2.E.34} $X,Y$ sets of sets, $X\in Y$; then $X\subseteq Y$.\ \textbf{2.E.35} $\varnothing\in\PS(X)$.\ \textbf{2.E.36} $\{\varnothing\}\in\PS(X)$.\ \textbf{2.E.37} $\{\varnothing,X\}\subseteq\PS(X)$.\ \textbf{2.E.39} $X\setminus Y=\varnothing$; then $X=Y$.\ \textbf{2.E.40} $X\setminus Y=Z$ and $Y\subseteq Z$; then $X=Y\cup Z$.\ \textbf{2.E.42} $A\subseteq X$; then $A\in X$.
\end{trybox}

